\section{Balanced Edge Scalars}\label{sec:peps_balanced_edge_scalars}

The edge gauges are determined only up to a scalar on each edge. The residual
freedom is the balanced edge scalars: an assignment of a nonzero scalar to every
oriented edge whose product around each vertex is one.
\begin{center}
    \TNPEPSVertexScalarBalance
\end{center}
Two gauge families related by such an assignment give the same gauge
equivalence, which is the precise sense in which the gauges of the Fundamental
Theorem are unique up to a constant.

\begin{definition}[Gauge equivalence modulo balanced edge scalars]%
    \label{def:GaugeEquivModEdgeScalars}
    \lean{TNLean.PEPS.GaugeEquivModEdgeScalars}
    \leanok
    \uses{def:edgeGaugeAt, def:edgeScalarAt, def:IsVertexBalanced}
    Two gauge families are equivalent modulo balanced edge scalars if there is
    a nonzero scalar $c_e$ on each edge such that the corresponding endpoint
    actions satisfy $X_{v,e}=c_{v,e}Y_{v,e}$, and
    $\prod_{e\ni v}c_{v,e}=1$ for every vertex $v$.
\end{definition}

\begin{theorem}[Multiplication of oriented edge scalars]%
    \label{thm:edgeScalarAt_mul}
    \lean{TNLean.PEPS.edgeScalarAt_mul}
    \leanok
    \uses{def:edgeScalarAt}
    If two edge-scalar families are multiplied, then their endpoint scalars
    satisfy $(cd)_{v,e}=c_{v,e}d_{v,e}$ on every oriented half-edge.
    \begin{center}
        \TNPEPSEdgeGaugeOrientation
    \end{center}
\end{theorem}

\begin{proof}\leanok
    At the lower endpoint this is ordinary multiplication of the edge scalars.
    At the upper endpoint, $(c_ed_e)^{-1}=c_e^{-1}d_e^{-1}$ because the
    edge scalars lie in $\mathbb C^\times$.
\end{proof}

\begin{theorem}[Inversion of oriented edge scalars]%
    \label{thm:edgeScalarAt_inv}
    \lean{TNLean.PEPS.edgeScalarAt_inv}
    \leanok
    \uses{def:edgeScalarAt}
    Inverting an edge-scalar family gives
    $(c^{-1})_{v,e}=c_{v,e}^{-1}$ on every oriented endpoint.
\end{theorem}

\begin{proof}\leanok
    At the lower endpoint, $(c^{-1})_{v,e}=c_e^{-1}=c_{v,e}^{-1}$.
    At the upper endpoint,
    $(c^{-1})_{v,e}=(c_e^{-1})^{-1}=c_{v,e}^{-1}$.
\end{proof}

\begin{theorem}[Nonzero oriented edge scalars]%
    \label{thm:edgeScalarAt_ne_zero}
    \lean{TNLean.PEPS.edgeScalarAt_ne_zero}
    \leanok
    \uses{def:edgeScalarAt}
    Every oriented endpoint scalar satisfies $c_{v,e}\ne 0$.
\end{theorem}

\begin{proof}\leanok
    Edge scalars are units, and the inverse of a unit is again a unit.
\end{proof}

\begin{theorem}[Product of balanced edge scalars]%
    \label{thm:IsVertexBalanced_mul}
    \lean{TNLean.PEPS.IsVertexBalanced.mul}
    \leanok
    \uses{def:IsVertexBalanced, thm:edgeScalarAt_mul}
    The pointwise product of two vertex-balanced edge-scalar families is again
    vertex-balanced.
\end{theorem}

\begin{proof}\leanok
    For every vertex $v$,
    $\prod_{e\ni v}(cd)_{v,e}=
        (\prod_{e\ni v}c_{v,e})(\prod_{e\ni v}d_{v,e})=1\cdot 1=1$.
\end{proof}

\begin{theorem}[Inverse of balanced edge scalars]%
    \label{thm:IsVertexBalanced_inv}
    \lean{TNLean.PEPS.IsVertexBalanced.inv}
    \leanok
    \uses{def:IsVertexBalanced, thm:edgeScalarAt_inv}
    The pointwise inverse of a vertex-balanced edge-scalar family is again
    vertex-balanced.
\end{theorem}

\begin{proof}\leanok
    For every vertex $v$,
    $\prod_{e\ni v}(c^{-1})_{v,e}=
        (\prod_{e\ni v}c_{v,e})^{-1}=1^{-1}=1$.
\end{proof}

\begin{theorem}[Reflexivity of balanced edge-scalar equivalence]%
    \label{thm:GaugeEquivModEdgeScalars_refl}
    \lean{TNLean.PEPS.GaugeEquivModEdgeScalars.refl}
    \leanok
    \uses{def:GaugeEquivModEdgeScalars}
    Every edge-gauge family is equivalent to itself modulo balanced edge
    scalars, witnessed by $c_e=1$ for every edge $e$.
    \begin{center}
        \TNPEPSEdgeGaugeOrientation
    \end{center}
\end{theorem}

\begin{proof}\leanok
    If $c_e=1$, then $c_{v,e}=1$, $X_{v,e}=c_{v,e}X_{v,e}$, and
    $\prod_{e\ni v}c_{v,e}=1$.
\end{proof}

\begin{theorem}[Symmetry of balanced edge-scalar equivalence]%
    \label{thm:GaugeEquivModEdgeScalars_symm}
    \lean{TNLean.PEPS.GaugeEquivModEdgeScalars.symm}
    \leanok
    \uses{def:GaugeEquivModEdgeScalars, thm:IsVertexBalanced_inv,
        thm:edgeScalarAt_ne_zero}
    If one gauge family differs from another by balanced edge scalars, then the
    second differs from the first by the inverse balanced edge-scalar family.
    \begin{center}
        \TNPEPSEdgeGaugeOrientation
    \end{center}
\end{theorem}

\begin{proof}\leanok
    If $X_{v,e}=c_{v,e}Y_{v,e}$ with $c_{v,e}\ne 0$, then
    $Y_{v,e}=c_{v,e}^{-1}X_{v,e}$. The inverse scalar family is balanced by
    Theorem~\ref{thm:IsVertexBalanced_inv}.
\end{proof}

\begin{theorem}[Transitivity of balanced edge-scalar equivalence]%
    \label{thm:GaugeEquivModEdgeScalars_trans}
    \lean{TNLean.PEPS.GaugeEquivModEdgeScalars.trans}
    \leanok
    \uses{def:GaugeEquivModEdgeScalars, thm:IsVertexBalanced_mul,
        thm:edgeScalarAt_mul}
    If one gauge family differs from a second by balanced edge scalars, and the
    second differs from a third by balanced edge scalars, then the first differs
    from the third by the pointwise product of the two scalar families.
    \begin{center}
        \TNPEPSEdgeGaugeOrientation
    \end{center}
\end{theorem}

\begin{proof}\leanok
    From $X_{v,e}=c_{v,e}Y_{v,e}$ and $Y_{v,e}=d_{v,e}Z_{v,e}$ one obtains
    $X_{v,e}=(c_{v,e}d_{v,e})Z_{v,e}$. The product scalar family is balanced
    by Theorem~\ref{thm:IsVertexBalanced_mul}.
\end{proof}

\begin{theorem}[The balanced edge-scalar quotient is an equivalence relation]%
    \label{thm:GaugeEquivModEdgeScalars_equivalence}
    \lean{TNLean.PEPS.GaugeEquivModEdgeScalars.equivalence}
    \leanok
    \uses{thm:GaugeEquivModEdgeScalars_refl,
        thm:GaugeEquivModEdgeScalars_symm,
        thm:GaugeEquivModEdgeScalars_trans}
    Gauge equivalence modulo balanced edge scalars is reflexive, symmetric,
    and transitive.
\end{theorem}

\begin{proof}\leanok
    Reflexivity is witnessed by $1$, symmetry by $c^{-1}$, and transitivity
    by $cd$.
\end{proof}

\begin{definition}[The balanced edge-scalar quotient relation]%
    \label{def:GaugeEquivModEdgeScalars_setoid}
    \lean{TNLean.PEPS.GaugeEquivModEdgeScalars.setoid}
    \leanok
    \uses{def:GaugeEquivModEdgeScalars,
        thm:GaugeEquivModEdgeScalars_equivalence}
    The preceding equivalence relation is the quotient relation on edge-gauge
    families modulo vertex-balanced edge scalars.
\end{definition}

\begin{theorem}[Balanced edge scalars preserve the local gauge action]%
    \label{thm:GaugeEquivModEdgeScalars_gaugeVertex_eq}
    \lean{TNLean.PEPS.GaugeEquivModEdgeScalars.gaugeVertex_eq}
    \leanok
    \uses{def:GaugeEquivModEdgeScalars, def:gaugeVertex}
    If two gauge families differ by a balanced edge-scalar reweighting, then
    they produce the same gauged tensor at every vertex.
\end{theorem}

\begin{proof}\leanok\uses{def:GaugeEquivModEdgeScalars, def:gaugeVertex}
    In each summand of the gauged tensor at $v$, the edge scalars contribute
    the factor $\prod_{e\ni v}c_{v,e}=1$. Hence the summand, and therefore
    the whole local tensor, is unchanged.
\end{proof}

\begin{theorem}[Balanced edge scalars preserve the gauged tensor]%
    \label{thm:GaugeEquivModEdgeScalars_applyGauge_eq}
    \lean{TNLean.PEPS.GaugeEquivModEdgeScalars.applyGauge_eq}
    \leanok
    \uses{thm:GaugeEquivModEdgeScalars_gaugeVertex_eq, def:applyGauge}
    If two gauge families differ by a balanced edge-scalar reweighting, then
    applying them to the same PEPS tensor produces the same gauged tensor.
    \begin{center}
        \TNPEPSGaugeVertexAction
    \end{center}
\end{theorem}

\begin{proof}\leanok
    For every vertex $v$, virtual index $\eta$, and physical index $\sigma$,
    $(A^X)_v(\eta,\sigma)=(A^Y)_v(\eta,\sigma)$ for the two gauged tensors by
    Theorem~\ref{thm:GaugeEquivModEdgeScalars_gaugeVertex_eq}. Hence the
    two gauged tensors are equal.
\end{proof}

\begin{theorem}[Balanced edge scalars preserve the gauged state]%
    \label{thm:GaugeEquivModEdgeScalars_applyGauge_sameState}
    \lean{TNLean.PEPS.GaugeEquivModEdgeScalars.applyGauge_sameState}
    \leanok
    \uses{thm:GaugeEquivModEdgeScalars_applyGauge_eq, def:peps_same_state}
    If two gauge families differ by a balanced edge-scalar reweighting, then
    the PEPS tensors obtained by applying those gauge families represent the
    same state.
\end{theorem}

\begin{proof}\leanok
    The two gauged tensors are equal, hence all of their state coefficients are
    equal.
\end{proof}

\begin{theorem}[Gauge uniqueness modulo balanced edge scalars]%
    \label{thm:gauge_unique_mod_edge_scalars}
    \lean{TNLean.PEPS.gauge_unique_mod_edge_scalars}
    \leanok
    \uses{def:GaugeEquivModEdgeScalars, def:IsVertexInjective,
        thm:peps_piProductFormsScalar}
    Let every bond space be nonzero. If two gauge families relate the same
    vertex-injective PEPS tensor to the same target tensor, then their edge
    gauges differ by nonzero scalars $c_e$ whose oriented product at every
    vertex is $1$. Thus the correct graph quotient is by balanced edge scalars.
    % Correction documented in docs/paper-gaps/peps_gauge_edge_scalars.tex.
    % The nonzero-bond restriction is the standing normal-PEPS convention,
    % matching the existence direction; see the same paper-gap note.
\end{theorem}

\begin{proof}\leanok
    Combining the two gauge relations gives, at every vertex, equality of the
    gauged local tensors. Linear independence of the local tensor family
    promotes this to equality of the products of incident edge-gauge entries
    over every pair of virtual configurations. Each oriented endpoint gauge is
    invertible and, since every bond space is nonzero, indexed by a nonempty
    type; the many-leg scalar extraction then yields, at each vertex, nonzero
    proportionality constants between the two families whose product is $1$.
    Reading the lower-endpoint constant of each edge defines a single edge
    scalar $c_e$; the inverse relation transports it to the upper endpoint, and
    invertibility makes the proportionality constant unique, so the vertex
    products of the per-vertex constants identify with the oriented products of
    $c_e$. These are all $1$, so $c$ is vertex-balanced.
\end{proof}

\begin{remark}
    The local left inverse $\Psi_v$, the realization map $T \mapsto O_T$,
    the split at one incident edge, and the partition
    $V = \{u\} \sqcup (V \setminus \{u,v\}) \sqcup \{v\}$ isolate the local
    indices and maps needed for the edge-centered reduction of
    \cite[Section~3]{Molnar2018NormalPEPS}. The middle-region blocked coefficient is
    constructed by the preceding definitions, and the local-gauge step is the
    comparison with the corresponding three-site injective chain. The corrected
    uniqueness statement is a quotient by balanced edge scalars.
\end{remark}
